% ============================================================
%  SparseRenderFormer — Eurographics / EGSR format
%  Compile with tectonic from this directory (all .cls/.sty present)
%  For submission: change \ConferencePaper to \ConferenceSubmission
%  and anonymise author block
% ============================================================

\documentclass{egpubl}
\usepackage{eg2026}

% --- publication type (one active at a time) ---
% \ConferenceSubmission   % anonymous submission
\ConferencePaper          % final camera-ready
% \SpecialIssuePaper      % CGF special issue (e.g. EGSR)

\CGFStandardLicense

% !! do not change anything above this line !!
% -----------------------------------------------------------
\usepackage[T1]{fontenc}
\usepackage{dfadobe}

\usepackage{cite}
\BibtexOrBiblatex
\electronicVersion
\PrintedOrElectronic

\ifpdf
  \usepackage[pdftex]{graphicx}
  \pdfcompresslevel=9
\else
  \usepackage[dvips]{graphicx}
\fi

\usepackage{egweblnk}
\usepackage{amsmath}
\usepackage{amssymb}
\usepackage{booktabs}
\usepackage{multirow}
\usepackage{xcolor}
\usepackage{microtype}
\usepackage{xspace}
% -----------------------------------------------------------

% Convenience macros
\newcommand{\ours}{SparseRenderFormer\xspace}
\newcommand{\RF}{RenderFormer\xspace}
\newcommand{\todo}[1]{\textcolor{red}{\textbf{[TODO: #1]}}}

\begin{document}

% ---------------------------------------------------------------------------
\title{Scaling Neural Global Illumination to Tens of Thousands of Triangles\\
       via Hierarchical Sparse Attention}

\author[V. Miene]{Victor Miene\\
  Carnegie Mellon University\\
  {\small\ttfamily vom@andrew.cmu.edu}
}

\maketitle

% ---------------------------------------------------------------------------
\begin{abstract}
Transformer-based methods for rendering resolve global illumination through
self-attention over triangle tokens.
The cost of this attention grows as $O(N^2)$ in the triangle count $N$,
so production meshes that routinely contain tens of thousands of triangles
can exceed practical memory and time limits.
In practice, most pairs of triangles exchange negligible radiance:
illumination is typically dominated by immediate neighbours, direct emitters,
and specular reflection paths.
Yet dense attention forces every triangle to attend to every other triangle,
paying the full $O(N^2)$ cost regardless of how little most pairs interact.
We present \ours, which exploits this structure by attending selectively
to a sparse subset of the triangles that matter most for each surface's
illumination, rather than all $N$.
This selective attention is organised into three branches structured around
the scene's bounding volume hierarchy: a coarse branch for global radiance,
a selected branch for direct illumination and specular paths, and a local
branch whose neighbourhood is restricted by the hemisphere constraint of the
rendering equation; both domain-specific constraints are physics-informed
and introduce no additional parameters.
\ours is initialised from the pretrained \RF checkpoint and fine-tuned
on progressively larger scenes.
At eight times \RF's 4{,}096-triangle training scale (32{,}768 triangles),
where dense attention exceeds practical memory limits, \ours achieves
\todo{XX}~dB PSNR with a wall-clock render time of \todo{XX}~ms per
frame versus \todo{XX}~ms for the dense baseline.
\end{abstract}

% ---------------------------------------------------------------------------
\section{Introduction}
\label{sec:intro}

Physically based rendering (PBR) simulates light transport by numerically
integrating the rendering equation~\cite{kajiya1986rendering}.
Production-quality path tracers such as Blender Cycles and Mitsuba~3~\cite{jakob2022drjit}
yield photorealistic results, but their cost grows with path depth and
sample count, which makes interactive or near-real-time global illumination
on complex scenes prohibitive.

Neural approximations have advanced along two lines.
Neural Radiance Fields (NeRF)~\cite{mildenhall2020nerf} and 3D Gaussian
Splatting (3DGS)~\cite{kerbl2023gaussian} achieve competitive view
synthesis but bake lighting into the scene representation, which limits
relighting.
Diffusion-based approaches such as DiffusionRenderer~\cite{liang2025diffusionrenderer}
achieve leading results on inverse and forward rendering of real video, but
operate on G-buffer representations that presuppose a rasterisation
pipeline and do not generalise across arbitrary mesh configurations.

\RF~\cite{zeng2025renderformer} introduces a different paradigm:
it encodes a triangle mesh scene as a sequence of tokens and resolves
global illumination end-to-end through transformer attention.
Each triangle token carries geometry (vertex positions), material (GGX
BRDF parameters~\cite{walter2007microfacet}), and emission.
A twelve-layer \emph{view-independent} stage performs triangle-to-triangle
self-attention to propagate inter-surface light transport.
A six-layer \emph{view-dependent} stage then lets ray-bundle tokens
cross-attend to the illuminated triangle sequence.
3D Rotary Position Encoding (3D~RoPE) follows the formulation
of~\cite{su2021roformer} and imprints geometric structure directly in
the attention logits.
The resulting 205M-parameter model is trained on 16M High Dynamic Range
(HDR) renders derived from Objaverse~\cite{deitke2022objaverse} and
generalises across materials, lighting, and geometry without per-scene
optimisation.

\paragraph{The scaling barrier.}
The view-independent stage grows as $O(N^2)$ in both memory and compute.
Doubling the triangle count quadruples the cost, so each step toward
production-scale meshes compounds the expense faster than any linear
increase in hardware capacity can offset.
Production assets average 8k--16k triangles after remeshing; architectural
scenes, characters, and product models routinely exceed 50k.
Reaching $N = 32{,}768$ from \RF's operating point of $N = 4{,}096$
requires $64\times$ more memory and compute under dense attention, a gap
that shifts the class of scenes that are renderable, not merely the cost.
\todo{Table: dense-attention PSNR and peak memory at $N \in \{4096, 8192, 16384\}$.}

\paragraph{Our approach.}
We exploit the observation that light transport is \emph{structurally
sparse}.
The rendering equation integrates radiance over the hemisphere, and most
distant surface pairs interact negligibly except through a small number of
high-albedo or specular paths.
This mirrors the structure that motivated Native Sparse Attention
(NSA)~\cite{yuan2025nsa} for language models: long-range context can be
captured at coarse resolution; a small fraction of highly relevant tokens
requires fine-grained attention; spatially adjacent tokens always deserve
local precision.

We formalise this structure in three branches, each grounded in the scene's BVH:
\begin{enumerate}
  \item \textbf{Coarse branch.} Pool triangle tokens within BVH leaf nodes
    by area-weighted averaging, then attend to compressed block tokens.
    Captures low-frequency global radiance.
  \item \textbf{Selected branch.} Each query triangle selects the
    top-$k$ triangle blocks by a learned relevance score.
    Captures direct illumination, one-bounce specular, and shadow casters.
  \item \textbf{Local branch.} Each query triangle attends to its
    $k_\ell$ spatially nearest triangles in 3D.
    Captures contact shadows and near-field colour bleeding.
\end{enumerate}
A per-head learned gate combines the three branch outputs.
Each triangle attends to $O(d \cdot b + k \cdot b + k_\ell)$ tokens,
where $b$ is the BVH block size, reducing the view-independent stage
from $O(N^2)$ to effectively $O(N)$ for fixed $d$, $k$, $k_\ell$.

\paragraph{Contributions.}
\begin{itemize}
  \item We introduce GP-Sparse, which replaces the $O(N^2)$
    view-independent self-attention in \RF with three BVH-structured
    attention branches and enables \ours to process scenes that exceed
    \RF's memory limit.
  \item Within the coarse branch, we introduce area-weighted pooling
    (Eq.~\ref{eq:coarse_pool}), in which each triangle's contribution
    to a compressed block token scales with its surface area, consistent
    with the rendering integral.
    The weights are precomputed at BVH build time and add no parameters.
  \item Within the local branch, we introduce normal-coherent masking
    (Eq.~\ref{eq:ncmask}), which suppresses back-facing neighbours by
    applying the hemisphere constraint of the rendering equation as a
    hard attention mask.
    The mask reduces the effective neighbourhood by approximately
    30--40\% on typical meshes and adds no inference cost.
  \item \todo{XX: PSNR / SSIM / memory results at 8k, 16k, 32k.}
\end{itemize}

% ---------------------------------------------------------------------------
\section{Related Work}
\label{sec:related}

\paragraph{Neural global illumination.}
Early methods cached irradiance in neural networks to accelerate Monte
Carlo sampling~\todo{cite}.
NeRF variants~\cite{mildenhall2020nerf} encode radiance as volumetric
implicit functions but conflate geometry and appearance.
3DGS~\cite{kerbl2023gaussian} achieves real-time synthesis via
differentiable point splatting but does not model inter-surface light
transport.
DiffusionRenderer~\cite{liang2025diffusionrenderer} couples a video
diffusion model with a deferred shading pipeline for inverse and forward
rendering, achieving leading results on real captures.
To our knowledge, \RF~\cite{zeng2025renderformer} is the first method to
model full global illumination over explicit triangle meshes via
end-to-end transformer attention without per-scene baking.

\paragraph{Efficient transformers.}
Full self-attention~\cite{vaswani2017attention} is $O(N^2)$ in both time and memory.
FlashAttention~\cite{dao2022flashattention} and its successors~\cite{dao2023flashattention2,shah2024flashattention3}
reduce memory by tiling the softmax computation but do not reduce
asymptotic complexity.
NSA~\cite{yuan2025nsa} jointly trains three hardware-aligned branches
(compressed token blocks, dynamically selected fine blocks, and a local
sliding window) and achieves near-lossless perplexity at $O(N)$ cost
on sequences up to 64k tokens.
The Differential Transformer~\cite{ye2024differential} sharpens
attention maps by cancelling noise through differential softmax pairs.
Our sparse attention design is closest to NSA but differs in block
construction (BVH rather than sequential blocks), positional encoding
(3D~RoPE applied within and across blocks rather than causal 1D~RoPE),
and neighbourhood construction: NSA and all prior sparse attention methods
operate on 1D token sequences with no spatial geometry, so their
neighbourhood masks carry no physical meaning.
To our knowledge, \ours is the first method to apply the hemisphere constraint of the
rendering equation as a hard geometric attention mask, implementing
back-face culling inside the local branch.

\paragraph{Equivariant scene encoders.}
SE(3)-equivariant transformers~\cite{liao2023equiformerv2} enforce
rotation and translation invariance by construction.
\RF achieves only translation invariance via 3D~RoPE and compensates
with rotation augmentation during training.
\ours inherits this limitation; equivariant attention is a complementary
future direction.

\paragraph{Hierarchical scene representations.}
Bounding Volume Hierarchies~\cite{kajiya1986rendering} organise primitives
into nested axis-aligned boxes for $O(\log N)$ ray intersection.
We repurpose BVH node boundaries as block boundaries for coarse pooling
and block selection, connecting classical rendering acceleration with
modern attention sparsity.

% ---------------------------------------------------------------------------
\section{Background: \RF Architecture}
\label{sec:background}

We briefly review the components of \RF~\cite{zeng2025renderformer} that
\ours modifies.

\paragraph{Scene tokenisation.}
A scene is represented as $N$ triangles.
Each triangle $i$ is encoded as a token $\mathbf{t}_i \in \mathbb{R}^D$
carrying its normal, GGX BRDF parameters $(\rho_d, \rho_s, \alpha)$, and
emission.
Geometric position is encoded via 3D~RoPE applied to attention logits,
not in the token itself.

\paragraph{3D Rotary Position Encoding.}
Given triangle $i$ with vertex positions
$\mathbf{v}_i^{(1)}, \mathbf{v}_i^{(2)}, \mathbf{v}_i^{(3)} \in \mathbb{R}^3$,
the RoPE bias between triangles $i$ and $j$ is:
\begin{equation}
  \phi_{ij} = \sum_{c \in \{1,2,3\}} \sum_{k=1}^{3}
  \operatorname{RoPE}\!\left(\mathbf{q}_i, \mathbf{v}_i^{(c)}\right)_k \cdot
  \operatorname{RoPE}\!\left(\mathbf{k}_j, \mathbf{v}_j^{(c)}\right)_k,
  \label{eq:rope3d}
\end{equation}
where $\mathbf{q}_i$ and $\mathbf{k}_j$ are the query and key projections
of the token embeddings.
This encodes relative 3D positions without global translation dependence.

\paragraph{View-independent stage.}
Twelve transformer layers perform full $N \times N$ self-attention over
triangle tokens to propagate inter-surface illumination.
This stage is the $O(N^2)$ bottleneck.

\paragraph{View-dependent stage.}
The camera view is discretised into $R$ ray-bundle tokens, each covering
an $8 \times 8$ pixel patch of the output image.
Six cross-attention layers produce HDR pixel-patch features from ray-bundle
queries attending to triangle keys and values, which are then decoded to colours.
Complexity is $O(R \cdot N)$ and becomes a secondary bottleneck at large $N$.

% ---------------------------------------------------------------------------
\section{\ours: BVH-Structured Sparse Attention}
\label{sec:method}

\subsection{Overview}

We replace the dense self-attention in the view-independent stage with a
three-branch sparse attention module.
All other components of \RF (tokenisation, 3D~RoPE, the view-dependent
stage, and the pixel decoder) are retained without modification.

\todo{Figure: architecture diagram.}

\subsection{BVH Block Construction}
\label{sec:bvh}

We build a BVH over the $N$ triangle centroids using surface area heuristic
(SAH) splitting.
At depth $d_\text{leaf}$, this yields $B = N / b$ leaf nodes, each
containing $b$ triangles on average.
We set $b = 32$ to match GPU warp width, which is consistent with
NSA's hardware-aligned design~\cite{yuan2025nsa}.
\todo{BVH library.}
Internal nodes at depth $d < d_\text{leaf}$ aggregate
$2^{d_\text{leaf} - d} \cdot b$ triangles.
The coarse branch pools at the parent level ($d_\text{leaf} - 1$), yielding
$B_p = N / (2b)$ parent blocks each covering $2b$ triangles on average.
The selected branch retrieves at the leaf level ($d_\text{leaf}$).
This two-level design gives the coarse branch a $2\times$ coarser global
view while the selected branch retains fine-grained retrieval, following
the two-granularity principle of NSA~\cite{yuan2025nsa}.

\subsection{Three-Branch Attention}
\label{sec:threebranch}

The two contributions of this paper modify the coarse and local branches
respectively; the selected branch is adopted from the NSA
design~\cite{yuan2025nsa} without modification.

Let $\mathbf{Q}, \mathbf{K}, \mathbf{V} \in \mathbb{R}^{N \times D}$ be
the query, key, and value matrices for a single attention head.

\paragraph{Coarse branch.}
Keys and values are pooled within each BVH leaf node by area-weighted
pooling, so that each triangle's contribution scales with its surface
area as it does under the rendering integral:
\begin{equation}
  \hat{\mathbf{K}}_\beta =
    \frac{\sum_{j \in \mathcal{B}_\beta} A_j\,\mathbf{K}_j}
         {\sum_{j \in \mathcal{B}_\beta} A_j},
  \qquad
  \hat{\mathbf{V}}_\beta =
    \frac{\sum_{j \in \mathcal{B}_\beta} A_j\,\mathbf{V}_j}
         {\sum_{j \in \mathcal{B}_\beta} A_j},
  \label{eq:coarse_pool}
\end{equation}
where $A_j = \tfrac{1}{2}\|(v_1^j{-}v_0^j)\times(v_2^j{-}v_0^j)\|$ is
the area of triangle $j$ and $\mathcal{B}_\beta$ is the triangle set of
leaf block $\beta$.
The weights are query-independent and are precomputed once at BVH build time.
To prevent triangles at leaf boundaries from being absent from neighbouring
blocks' pooled tokens, the $k_\text{fringe} = 4$ geometrically closest
triangles from adjacent leaves are included in each leaf's pooling with
their true areas.
Fringe membership is determined at BVH build time and adds no per-forward-pass
overhead.
Each query $\mathbf{Q}_i$ then attends to all $B_p = N/(2b)$ parent-level
compressed block tokens:
\begin{equation}
  \mathbf{O}_i^\text{coarse} =
    \operatorname{Attn}(\mathbf{Q}_i,\, \hat{\mathbf{K}}^p,\, \hat{\mathbf{V}}^p),
\end{equation}
where $\hat{\mathbf{K}}^p, \hat{\mathbf{V}}^p$ are parent-level pooled tokens.
This branch captures low-frequency radiance transfer and runs in
$O(N \cdot B_p) = O(N^2 / (2b))$, a $64\times$ reduction in keys per
query at $b = 32$.

\paragraph{Selected branch.}
For each query triangle $i$, the top-$k_\text{sel}$ leaf-level blocks are
retrieved using a dedicated lightweight block scorer with $H^I$ heads:
\begin{equation}
  I_{i,\beta} = \sum_{j=1}^{H^I} w^I_{i,j}
                \cdot \mathrm{ReLU}\!\left(q^I_{i,j} \cdot \bar{k}^I_\beta\right),
  \label{eq:scorer}
\end{equation}
where $q^I_{i,j} \in \mathbb{R}^{d^I}$ and $w^I_{i,j} \in \mathbb{R}$ are
derived from query token $\mathbf{Q}_i$, and $\bar{k}^I_\beta \in \mathbb{R}^{d^I}$
is derived from the area-weighted pooled token of leaf block $\beta$
(Eq.~\ref{eq:coarse_pool}).
ReLU activation suppresses negative scores without normalisation, which is
more hardware-efficient than softmax for large block counts.
The discrete top-$k$ selection is made end-to-end differentiable via a
STE~\cite{bengio2013estimating}: the forward pass performs hard top-$k$
retrieval, and the backward pass treats the scorer logits $I_{i,:}$ as
continuous weights, passing gradients straight through the argmax into
the scorer and, via $q^I_{i,j}$, into the query representations of the
main model.
An auxiliary KL divergence objective (Section~\ref{sec:finetune}) provides
a calibration signal during Stage~1 warm-up.
Full-resolution leaf tokens for the top-$k_\text{sel}$ blocks are gathered:
\begin{equation}
  \mathcal{S}_i = \operatorname{top-}k_\text{sel}\!\left\{ I_{i,\beta} \right\}_\beta,
  \quad
  \mathbf{O}_i^\text{sel} =
    \operatorname{Attn}\!\left(
      \mathbf{Q}_i,\; \mathbf{K}_{\mathcal{S}_i},\; \mathbf{V}_{\mathcal{S}_i}
    \right).
  \label{eq:sel}
\end{equation}
Each query attends to $k_\text{sel} \cdot b$ fine-grained tokens, targeting
direct illumination and specular paths where precise geometry determines the
outcome.
The default settings are $k_\text{sel} = 4$ and $H^I = 4$.

\paragraph{Local branch.}
Each query triangle $i$ attends to its $k_\ell = 64$ spatially nearest
triangles in 3D.
On a typical mesh, however, 30--40\% of the nearest neighbours in
Euclidean space face away from the query triangle.
Under the rendering equation, a triangle whose normal points into the
lower hemisphere of triangle $i$ contributes zero radiance, because the
cosine term $(\mathbf{n}_i \cdot \omega)$ clamps at zero for directions
below the surface.
Attending to these neighbours wastes capacity and degrades shadow
boundary quality.

We apply a \emph{normal-coherent mask} that suppresses back-facing
neighbours before the softmax, implementing the transformer equivalent
of back-face culling.
After computing $\mathcal{L}_i$, each neighbour $j$ is masked if its
normal exceeds a $120^\circ$ angle from the query normal.
The mask is defined by a sigmoid soft approximation and its hard
inference counterpart:
\begin{align}
  \tilde{m}_{ij} &= \sigma\!\bigl(-s\,(\mathbf{n}_i \cdot \mathbf{n}_j + \tau)\bigr),
  \label{eq:ncmask_soft} \\
  m_{ij}         &= \mathbf{1}[\mathbf{n}_i \cdot \mathbf{n}_j < -\tau],
  \label{eq:ncmask_hard}
\end{align}
where $\tau = 0.5$ and sharpness $s = 20$.
The threshold $\tau = 0.5$ encodes the hemisphere constraint of the
rendering integral, not a learned preference.
To maintain end-to-end differentiability during training we apply a
\emph{straight-through estimator} (STE)~\cite{bengio2013estimating}: the
forward pass uses the hard indicator $m_{ij}$ (physically exact), while
the backward pass propagates gradients through the soft sigmoid
$\tilde{m}_{ij}$, so the normal vectors $\mathbf{n}_i$ and $\mathbf{n}_j$
receive gradient signal from the rendering loss.
At inference only $m_{ij}$ is evaluated.
The masked local attention is then:
\begin{align}
  \mathcal{L}_i &= \operatorname{kNN}_{k_\ell}\!(c_i,\, \{c_j\}_{j \neq i}),
  \label{eq:knn} \\
  \mathbf{O}_i^\text{local} &=
    \operatorname{softmax}\!\!\left(
      \frac{\mathbf{Q}_i\,\mathbf{K}_{\mathcal{L}_i}^\top}{\sqrt{D}}
      - \mathbf{m}_i^{\mathrm{STE}} \cdot \infty
    \right)
    \mathbf{V}_{\mathcal{L}_i},
  \label{eq:local_masked}
\end{align}
where $\mathbf{m}_i^{\mathrm{STE}} \in \mathbb{R}^{k_\ell}$ uses $m_{ij}$
in the forward pass and $\tilde{m}_{ij}$ in the backward pass.
This branch captures contact shadows and near-field colour bleeding.
The mask reduces the effective neighbourhood by approximately 30--40\%
on Objaverse scenes, at no added inference cost.
\todo{kNN algorithm.}
\todo{Attention kernel.}

\paragraph{Gated combination.}
Per-head learned scalar gates $\lambda_1, \lambda_2, \lambda_3 \geq 0$
(softmax-normalised) combine the three branch outputs:
\begin{equation}
  \mathbf{O}_i =
    \lambda_1 \mathbf{O}_i^\text{coarse}
    + \lambda_2 \mathbf{O}_i^\text{sel}
    + \lambda_3 \mathbf{O}_i^\text{local}.
  \label{eq:gate}
\end{equation}

\subsection{3D RoPE in Sparse Branches}
\label{sec:rope}

In the coarse branch, the position argument for compressed block token
$\hat{\mathbf{K}}_\beta$ is the centroid of all triangles in
$\mathcal{B}_\beta$.
In the selected and local branches, each token retains its original
per-triangle 3D position as in Eq.~\ref{eq:rope3d}.
The coarse branch therefore encodes block-level geometry while the fine
branches encode triangle-level geometry.

\subsection{Sparse View-Dependent Cross-Attention}
\label{sec:viewdep}

At large $N$, the $O(R \cdot N)$ cross-attention from ray bundles to
triangles also becomes a bottleneck.
Frustum culling and depth-sorted $k$-NN selection limit each ray-bundle
token to attending to the $k_r = 256$ triangles most likely to be visible
from that ray's frustum, reducing view-dependent complexity to
$O(R \cdot k_r)$ independent of $N$.

\subsection{Fine-Tuning from Pretrained \RF Weights}
\label{sec:finetune}

\RF's pretrained 205M-parameter checkpoint provides initialisation.
All weights transfer except the attention projection matrices in the
view-independent stage, which are re-initialised.
Fine-tuning runs for \todo{XX epochs} on a filtered subset of the \RF
training distribution~\cite{deitke2022objaverse} with $N \leq 32{,}768$
triangles, using Blender Cycles renders as pseudo-ground-truth.
We use the \todo{optimiser} with an initial learning rate
of \todo{XX} and a cosine decay schedule following a \todo{XX}-step
linear warmup.
The loss is \todo{loss function}.
Batch size is \todo{XX}.
Mixed-precision training uses \todo{mixed precision settings}.
Fine-tuning runs on \todo{hardware and training time}.
The Swin-Large backbone~\cite{liu2021swin} in the view-dependent pixel
decoder is frozen throughout.

\paragraph{Two-stage training procedure.}
Fine-tuning proceeds in two stages, following the indexer warm-up principle
introduced in DeepSeek Sparse Attention~\cite{deepseek2025dsa}.

\textbf{Stage 1: Dense warm-up.}
All GP-Sparse parameters are frozen except the block scorer
(Eq.~\ref{eq:scorer}).
Full dense attention runs normally.
The scorer is trained with a KL divergence loss against the L1-normalised
dense attention distribution:
\begin{equation}
  \mathcal{L}^I = \sum_i
    \mathbb{D}_{\mathrm{KL}}\!\left(
      p_{i,:} \;\Big\|\; \mathrm{Softmax}(I_{i,:})
    \right),
  \label{eq:scorer_kl}
\end{equation}
where $p_{i,:}$ is the sum of the dense attention weights for triangle $i$
across all heads, L1-normalised over the sequence dimension.
This stage trains the scorer to predict which BVH blocks the dense model
attends to before any sparse routing is imposed.
The warm-up runs for \todo{XX} steps on the same data distribution as Stage~2.

\textbf{Stage 2: Sparse fine-tuning.}
Sparse block selection and the normal-coherent mask are activated; all
parameters are unfrozen.
Both discrete operations use their respective STEs
(Section~\ref{sec:threebranch}), so gradients from the rendering loss flow
end-to-end through the selected branch into the block scorer and through
the local branch into the triangle normal representations.
The auxiliary $\mathcal{L}^I$ continues as a secondary calibration signal
for the scorer, weighted by $\alpha = 0.01$:
$\mathcal{L} = \mathcal{L}^{\text{render}} + \alpha\,\mathcal{L}^I$.

\paragraph{Complexity summary.}
\begin{center}\small
\begin{tabular}{lcc}
\toprule
Stage & \RF & \ours \\
\midrule
View-indep. & $O(N^2)$ & $O(N(N/b + k_s b + k_\ell))$ \\
View-dep. & $O(RN)$ & $O(Rk_r)$ \\
\bottomrule
\end{tabular}
\end{center}
For $N = 32{,}768$, $b = 32$, $k_\text{sel} = 4$, $k_\ell = 64$,
$k_r = 256$, \ours attends to approximately \todo{XX$\times$} fewer token
pairs than \RF extrapolated to the same $N$.

% ---------------------------------------------------------------------------
\section{Experiments}
\label{sec:experiments}

\todo{Complete after implementation.}

\subsection{Setup}

\paragraph{Scenes.}
Three scene categories form the evaluation suite:
(1) synthetic Cornell-box variants (100 scenes, $N \leq 2{,}048$)
for sanity checking;
(2) remeshed Objaverse scenes~\cite{deitke2022objaverse}
(1{,}000 scenes, $2{,}048 \leq N \leq 32{,}768$) for main comparison;
(3) complex single-room scenes (50 scenes,
$32{,}768 \leq N \leq 65{,}536$) to probe out-of-distribution triangle counts.
Objaverse meshes are remeshed using a signed-distance-field
marching-cubes reconstruction followed by Qslim face decimation
to obtain triangle counts in the target range.
The fine-tuning split comprises 10{,}000 scenes rendered from 4
viewpoints each, using the same random-rotation, random-material, and
random-lighting augmentation as \RF~\cite{zeng2025renderformer}.
Pseudo-ground-truth images are produced with Blender Cycles at 4{,}096~spp
and $512 \times 512$ resolution.

\paragraph{Baselines.}
We compare against three baselines:
(1) the released \RF checkpoint with dense attention, evaluated at
$N \leq 4{,}096$;
(2) \RF fine-tuned with dense attention up to the largest $N$ at which
$O(N^2)$ attention remains tractable ($N \leq 8{,}192$);
(3) Blender Cycles reference renders (512~spp, $\approx$90\,s per frame)
used as ground truth.

\paragraph{Metrics.}
Peak Signal-to-Noise Ratio (PSNR), Structural Similarity (SSIM), and
LPIPS~\cite{zhang2018lpips} are computed against Blender Cycles reference
renders at $512 \times 512$ resolution.
Peak GPU memory (GB) and wall-clock render time per frame (ms) are
measured on a single A100-80GB.

% ---------------------------------------------------------------------------
\subsection{Main Results}
\label{sec:results}

Table~\ref{tab:main} reports PSNR, SSIM, and LPIPS at
$N \in \{4{,}096,\, 8{,}192,\, 16{,}384,\, 32{,}768\}$.
At $N = 32{,}768$, \ours achieves \todo{XX}~dB PSNR,
\todo{XX} SSIM, and \todo{XX} LPIPS.
The dense-attention baseline fine-tuned to $N = 8{,}192$ reaches
\todo{XX}~dB PSNR and cannot be evaluated beyond that scale.
Peak GPU memory is reduced from \todo{XX}~GB (dense, extrapolated to
$N = 32{,}768$) to \todo{XX}~GB, and wall-clock render time per frame
decreases from \todo{XX}~ms to \todo{XX}~ms.

\begin{table}[h]
  \caption{PSNR\,/\,SSIM\,/\,LPIPS vs baselines. OOM denotes
           out-of-memory on a single A100-80GB.}
  \label{tab:main}
  \centering
  \resizebox{\columnwidth}{!}{%
  \begin{tabular}{lrccc}
    \toprule
    Method & $N$ & PSNR$\uparrow$ & SSIM$\uparrow$ & LPIPS$\downarrow$ \\
    \midrule
    \RF (pretrained)  & 4{,}096  & \todo{XX} & \todo{XX} & \todo{XX} \\
    \RF (fine-tuned)  & 8{,}192  & \todo{XX} & \todo{XX} & \todo{XX} \\
    \RF (fine-tuned)  & 16{,}384 & OOM       & $-$       & $-$       \\
    \midrule
    \ours             & 8{,}192  & \todo{XX} & \todo{XX} & \todo{XX} \\
    \ours             & 16{,}384 & \todo{XX} & \todo{XX} & \todo{XX} \\
    \ours             & 32{,}768 & \todo{XX} & \todo{XX} & \todo{XX} \\
    \bottomrule
  \end{tabular}%
  }
\end{table}

Figure~\ref{fig:scaling} shows peak GPU memory and PSNR as a function of
$N$ on a log-log scale, confirming the transition from $O(N^2)$ to
$O(N)$ scaling.
\todo{Figure~\ref{fig:scaling}: PSNR and peak memory vs $N$.}

Figure~\ref{fig:qual} shows side-by-side renders of \ours, \RF, and the
Blender Cycles reference on three representative scenes, including
$5\times$ difference images.
\todo{Figure~\ref{fig:qual}: qualitative comparison.}

Figure~\ref{fig:fail} illustrates two failure modes specific to GP-Sparse:
a concave-geometry scene where the hemisphere-mask fallback fires, and a
translucent-material scene where the mask incorrectly suppresses
back-hemisphere neighbours.
\todo{Figure~\ref{fig:fail}: failure cases.}

% ---------------------------------------------------------------------------
\subsection{Ablation Studies}
\label{sec:ablation}

\paragraph{Branch ablation.}
Each branch is ablated by setting its gate weight $\lambda = 0$ during
inference (Eq.~\ref{eq:gate}).
Removing the coarse branch reduces PSNR by \todo{XX}~dB;
removing the selected branch reduces PSNR by \todo{XX}~dB;
removing the local branch reduces PSNR by \todo{XX}~dB.
The full three-branch model outperforms the best single-branch ablation
by \todo{XX}~dB at $N = 16{,}384$.

\begin{table}[h]
  \caption{Branch ablation at $N = 16{,}384$. \checkmark~denotes branch enabled.}
  \label{tab:branch}
  \centering\small
  \begin{tabular}{ccccc}
    \toprule
    Coarse & Selected & Local & PSNR$\uparrow$ & SSIM$\uparrow$ \\
    \midrule
    \checkmark & \checkmark & \checkmark & \todo{XX} & \todo{XX} \\
               & \checkmark & \checkmark & \todo{XX} & \todo{XX} \\
    \checkmark &            & \checkmark & \todo{XX} & \todo{XX} \\
    \checkmark & \checkmark &            & \todo{XX} & \todo{XX} \\
    \bottomrule
  \end{tabular}
\end{table}

\paragraph{Block size and selection count.}
At $b = 32$ and $k_\text{sel} = 4$, \ours reaches \todo{XX}~dB PSNR at
\todo{XX}~GB peak memory.
Increasing $b$ to 64 reduces memory by \todo{XX}\% but degrades PSNR by
\todo{XX}~dB; reducing $k_\text{sel}$ to 2 has a \todo{smaller / larger}
effect on quality than on memory.

\begin{table}[h]
  \caption{PSNR and peak memory vs block size $b$ and selection count
           $k_\text{sel}$ at $N = 16{,}384$.}
  \label{tab:blocksize}
  \centering\small
  \begin{tabular}{cccc}
    \toprule
    $b$ & $k_\text{sel}$ & PSNR$\uparrow$ & Mem (GB) \\
    \midrule
    16 & 4 & \todo{XX} & \todo{XX} \\
    32 & 4 & \todo{XX} & \todo{XX} \\
    64 & 4 & \todo{XX} & \todo{XX} \\
    32 & 2 & \todo{XX} & \todo{XX} \\
    32 & 8 & \todo{XX} & \todo{XX} \\
    \bottomrule
  \end{tabular}
\end{table}

\paragraph{Position encoding in coarse branch.}
Block centroid RoPE outperforms no positional encoding by \todo{XX}~dB
and per-triangle RoPE applied to pooled keys by \todo{XX}~dB at
$N = 16{,}384$, confirming that block-level geometry is sufficient for
the coarse branch.

% ---------------------------------------------------------------------------
\subsection{Scaling}
\label{sec:scaling}

Figure~\ref{fig:scaling} shows that dense attention memory grows as
$O(N^2)$, reaching \todo{XX}~GB at $N = 16{,}384$ and becoming infeasible
beyond $N = \todo{XX}$.
\ours scales approximately linearly, using \todo{XX}~GB at
$N = 32{,}768$ and \todo{XX}~GB at $N = 65{,}536$.
Render time follows the same trend: the dense baseline requires
\todo{XX}~ms per frame at $N = 8{,}192$, while \ours requires
\todo{XX}~ms at $N = 32{,}768$.
\todo{Figure~\ref{fig:scaling}: log-log memory and time vs $N$.}

% ---------------------------------------------------------------------------
\section{Conclusion}
\label{sec:conclusion}

We have presented \ours. The architecture scales transformer-based neural
global illumination from 4{,}096 to over 32{,}768 triangles per scene by
replacing dense self-attention with a BVH-structured three-branch sparse
module.
Each branch targets a distinct scale of light transport: the coarse branch
resolves global radiance, the selected branch targets specular and shadow
paths, and the local branch handles near-field contact effects.
Fine-tuning from the released \RF checkpoint makes the approach tractable
without an 8-day multi-GPU training run.

\paragraph{Limitations and future work.}
\ours inherits \RF's restriction to white diffuse lights outside the scene
bounding box and GGX BRDFs without spatially varying textures.
\todo{GP-Sparse limitations: concave geometry, translucent materials, noisy normals, query-independent area weights.}
Extensions to environment lighting, HDRI inputs, spatially varying BRDFs,
and SO(3)-equivariant attention~\cite{liao2023equiformerv2} remain open
directions.
End-to-end differentiability via the STEs in Section~\ref{sec:threebranch}
enables inverse rendering applications, which we plan to explore.

% ---------------------------------------------------------------------------
\bibliographystyle{eg-alpha-doi}
\bibliography{references}

\end{document}
